% ═══════════════════════════════════════════════════════════════════════════
%  PAPER SECTION: AVMI UGV Dataset, Model, Cross-Dataset Transfer
%  Drop this into the main paper where indicated.
% ═══════════════════════════════════════════════════════════════════════════

% ───────────────────────────────────────────────────────────────────────────
\section{AVMI UGV Dataset and Terrain Taxonomy}
\label{sec:dataset}

\subsection{Terrain Class Ontology}

Existing off-road datasets define broad class ontologies intended for general
environment description rather than navigation-specific perception.  RUGD~\cite{rugd2019}
provides 24 classes, and RELLIS-3D~\cite{rellis2021} provides 20 classes, yet both
suffer from severe pixel-level class imbalance: in RUGD, grass and terrain
together account for more than 80\% of all annotated pixels, while classes
such as \textit{bicycle} or \textit{sign} contribute less than 1\%.
RELLIS-3D similarly concentrates the label mass in sky, grass, tree, and
bush ($\sim$94\% of labeled pixels), leaving critical navigation-relevant
classes under-represented.

Our AVMI UGV dataset addresses these limitations by adopting a compact,
\emph{visually grounded} six-class taxonomy designed around what a wheeled
ground vehicle must perceive to navigate safely in outdoor terrain:

\begin{itemize}
    \item \textbf{Sky} --- open sky and clouds.  Always non-traversable;
          provides strong spatial context for horizon detection.
    \item \textbf{Tree} --- tree trunks and full canopy.  Non-traversable;
          key landmark in forest and trail environments.
    \item \textbf{Bush} --- low shrubs and dense ground-level vegetation.
          May be partially traversable but is treated as an obstacle for
          conservative planning.
    \item \textbf{Ground} --- open grass, bare earth, and low-friction
          surfaces.  The primary traversable class.
    \item \textbf{Rock} --- exposed rock faces and loose rock beds.
          Traversable only with caution; associated with high terrain
          roughness.
    \item \textbf{Stump} --- wooden stumps, logs, and similar low-lying
          solid obstacles.  Non-traversable and frequently missed by
          navigability-focused ontologies.
\end{itemize}

Table~\ref{tab:ontology} compares our ontology with RUGD and RELLIS-3D and
shows how each external class maps to our six-class scheme.  Classes that
lack a visually consistent counterpart in our taxonomy (e.g., asphalt,
concrete, fence) are either merged conservatively or left unmapped depending
on the transfer strategy (Section~\ref{sec:mapping}).

\begin{table}[h]
\centering
\caption{Comparison of dataset class ontologies and AVMI mapping.
         ``---'' indicates no mapping (ignored during cross-dataset training).}
\label{tab:ontology}
\renewcommand{\arraystretch}{1.15}
\begin{tabular}{p{2cm}p{2.2cm}p{2.2cm}}
\toprule
\textbf{AVMI class} & \textbf{RUGD sources} & \textbf{RELLIS-3D sources} \\
\midrule
Sky      & sky                   & sky                \\
Tree     & tree                  & tree               \\
Bush     & bush                  & bush               \\
Ground   & dirt, gravel,         & dirt, mud,         \\
         & (grass$^\dagger$)     & (grass$^\dagger$)  \\
Rock     & rock, rock-bed,       & rubble             \\
         & water$^\ddagger$      &                    \\
Stump    & log, (pole)           & log                \\
\midrule
\textit{Ignored} & sand, asphalt, concrete, & grass, asphalt, concrete, \\
         & mulch, vehicle, building, & pole, water, vehicle, \\
         & fence, sign, bridge,  & building, fence, barrier, \\
         & picnic-table, bicycle & puddle, person     \\
\bottomrule
\multicolumn{3}{l}{\small $^\dagger$ Grass is visually green like tree/bush; mapping to Ground}\\
\multicolumn{3}{l}{\small causes visual contradiction --- see Section~\ref{sec:mapping}.}\\
\multicolumn{3}{l}{\small $^\ddagger$ Water mapped conservatively as non-traversable (Rock class).}
\end{tabular}
\end{table}

\subsection{Dataset Collection and Annotation}

The AVMI UGV dataset is collected using Unreal Engine (UE) simulation,
rendered from the perspective of a ground vehicle traversing open grassy
terrain interspersed with trees, rocks, and stumps.  This simulation-based
approach allows us to generate pixel-perfect annotations without manual
labelling effort, while controlling the visual diversity (lighting conditions,
time of day, vegetation density) and ensuring a more balanced class
distribution than is achievable with real-world datasets alone.

Each frame is annotated at full pixel precision using UE's semantic
segmentation pass, which assigns a unique RGB colour to each object class.
At training time, annotations are converted to integer class indices via
nearest-colour L2 matching in RGB space:
\begin{equation}
    c^* = \operatorname*{arg\,min}_{c \in \mathcal{C}} \| \mathbf{p} - \mathbf{r}_c \|_2
\end{equation}
where $\mathbf{p}$ is the RGB value of an annotation pixel and $\mathbf{r}_c$
is the reference colour for class $c$.  This matching step is necessary to
correct minor JPEG compression artefacts that shift annotation colours by a
few intensity values.

Frames with severe motion blur or insufficient class coverage (e.g., frames
where fewer than two classes are present) are excluded from the dataset.
The retained frames are split into disjoint training, validation, and
held-out test sets.  Figure~\ref{fig:avmi_sample} shows representative
frames from the dataset with their ground-truth annotation overlays.

\begin{figure}[h]
    \centering
    \includegraphics[width=\linewidth]{results/avmi_test_batch/avmi_sample.png}
    \caption{Representative AVMI UGV frames.  Left: RGB image.
             Centre: ground-truth segmentation mask
             (blue=sky, green=tree, yellow=bush, brown=ground,
             red=rock, pink=stump).
             Right: overlay.}
    \label{fig:avmi_sample}
\end{figure}

% ───────────────────────────────────────────────────────────────────────────
\section{Model Architecture and Training}
\label{sec:model}

\subsection{Network Architecture}

We adopt GANav~\cite{ganav2022}, an attention-based semantic segmentation
network designed for texture-driven off-road terrain boundaries.  The key
components are:

\begin{itemize}
    \item \textbf{Backbone} --- MixVisionTransformer-B0 (MiT-B0) with four
          hierarchical stages, embedding dimension 32, and spatial reduction
          ratios $[8,4,2,1]$.  Outputs multi-scale feature maps at
          $1/8$, $1/16$, $1/32$, and $1/64$ of the input resolution.
    \item \textbf{Decode head} --- \texttt{OursHeadClassAtt}, which fuses
          the four-scale features and applies a bi-directional Polarised
          Self-Attention (PSA) module over a fixed $97\times97$ attention
          mask.  The fused feature channel dimension is 384.
    \item \textbf{Auxiliary heads} --- Two lightweight FCN heads at the
          second and third backbone stages provide intermediate deep
          supervision during training (loss weights 0.4 each).
\end{itemize}

Input images are randomly cropped to $300\times375$ pixels and normalised
with ImageNet statistics ($\mu=[123.7, 116.3, 103.5]$,
$\sigma=[58.4, 57.1, 57.4]$).  The PSA attention mask size of $97\times97$
is derived from the $300\times375$ crop and must be preserved across all
experiments; changing the crop size invalidates the learned attention masks.
The total model size is approximately $3.7\,\text{M}$ parameters.

All experiments are run on an NVIDIA RTX 5070 Laptop GPU (CUDA compute
capability sm\_120).  Because PyTorch's cuDNN does not yet support sm\_120,
cuDNN is disabled (\texttt{torch.backends.cudnn.enabled = False}) and
inference falls back to native CUDA kernels.

\subsection{Training from Scratch on the AVMI Dataset}

The model is trained from scratch on the AVMI UGV dataset for
\textbf{240,000 iterations} using stochastic gradient descent:
\begin{itemize}
    \item Learning rate: $0.06$ with polynomial decay (power $0.9$,
          minimum $10^{-4}$)
    \item Linear warmup for 1,500 iterations (initial ratio $10^{-6}$)
    \item Momentum: $0.9$; weight decay: $4\times10^{-5}$
    \item Batch size: $4$ (2 images per GPU)
    \item No external pretrained weights (ImageNet or otherwise)
\end{itemize}

Table~\ref{tab:scratch} reports per-class and mean IoU on the AVMI
validation set at convergence.  Sky and ground are well-learned from the
start due to their visual distinctiveness and large pixel count.  Tree is
also well-segmented (IoU $70.2\%$) because the simulation renders tree
canopy with a consistent dark-green texture.  Bush and stump remain
challenging: bush boundaries are diffuse and visually similar to ground
vegetation, while stump pixels are sparse across the dataset.

\begin{table}[h]
\centering
\caption{AVMI scratch model --- per-class IoU (\%) on validation set
         at 240k iterations.}
\label{tab:scratch}
\begin{tabular}{lcccccc|c}
\toprule
 & Sky & Tree & Bush & Ground & Rock & Stump & \textbf{mIoU} \\
\midrule
IoU & 86.9 & 70.2 & 4.9 & 91.2 & 40.4 & 15.5 & \textbf{51.51} \\
\bottomrule
\end{tabular}
\end{table}

This trained model provides an AVMI-specific feature space that serves as
the initialisation for all cross-dataset fine-tuning experiments described
below.  We refer to it as the \emph{base model} throughout.

\subsection{Cross-Dataset Class Mapping Strategies}
\label{sec:mapping}

To evaluate how well the AVMI-trained feature space generalises to real
outdoor data, we fine-tune the base model on RUGD and RELLIS-3D.
Because these datasets use different class ontologies, each annotation must
be remapped to our six-class scheme before training.  We investigated three
strategies, each revealing a distinct failure mode.

\subsubsection{Full Mapping}

All source classes are assigned to the nearest AVMI class.  For RUGD, this
means dirt, sand, grass, asphalt, gravel, mulch, and concrete are all
merged into \textit{Ground}; pole, vehicle, fence, building, log, sign, and
bridge are merged into \textit{Stump} (obstacle); rock and rock-bed into
\textit{Rock}.

\textbf{Problem --- class imbalance.}  Under full mapping, \textit{Ground}
accounts for \textbf{68\%} of all RUGD labeled pixels (Table~\ref{tab:dist}).
Because cross-entropy loss is proportional to pixel count, the gradient is
dominated by Ground, and the network stops learning to distinguish trees or
sky.  After 100k iterations, IoU.tree remained at only $5.6\%$ and
IoU.bush at $0.0\%$ despite the base model achieving IoU.tree $= 70.2\%$.
The model had simply forgotten the tree class under Ground's overwhelming
gradient pressure.

\begin{table}[h]
\centering
\caption{Effective pixel distribution (\%) among labeled pixels after each
         RUGD mapping strategy.  ``n/a'' means the class has no pixels
         (obstacle not mapped in selective).}
\label{tab:dist}
\begin{tabular}{lccc}
\toprule
\textbf{Class} & \textbf{Full map} & \textbf{Selective} & \textbf{Full+Weight} \\
\midrule
Sky      &  2.0 & 25.5 &  2.0 \\
Tree     &  0.2 &  3.7 &  0.2 \\
Bush     &  0.0 &  0.0 &  0.0 \\
Ground   & 67.8 & 52.8 & 67.8 \\
Rock     & 12.3 & 18.1 & 12.3 \\
Stump    & 17.6 &  n/a & 17.6 \\
\bottomrule
\end{tabular}
\end{table}

\subsubsection{Selective Mapping}

To prevent Ground from overwhelming the loss, we mapped \emph{only} the
classes whose visual appearance is directly consistent with an AVMI class,
and set all other pixels to the ignore index (255):
\begin{itemize}
    \item RUGD: dirt$\to$Ground, gravel$\to$Ground, tree$\to$Tree,
          sky$\to$Sky, rock/rock-bed$\to$Rock
    \item RELLIS: dirt$\to$Ground, mud$\to$Ground, tree$\to$Tree,
          sky$\to$Sky, bush$\to$Bush, rubble$\to$Rock
    \item Everything else $\to$ \textbf{ignore (255)}
\end{itemize}

This redistributes the labelled pixel budget: sky rises from 2.0\% to
25.5\% and tree from 0.2\% to 3.7\% of actively trained pixels
(Table~\ref{tab:dist}).  On the RUGD validation set, this yielded the
best RUGD mIoU of $67.4\%$ with IoU.tree $= 58.7\%$ and IoU.sky $= 87.4\%$.

\textbf{Problem --- coverage gap.}  With 94\% of RUGD pixels ignored,
the model has never received supervision for grass, asphalt, or concrete
textures.  On unseen UGV images (which contain green grass ground), the
model produces random class predictions for these textures, resulting in
visually incoherent outputs.  The selective strategy works well within its
source domain but fails to generalise to the UGV test domain.

\subsubsection{RELLIS Index Mapping Correction}

The RELLIS-3D annotation files (\texttt{\_orig.png}) store
\emph{sequential} indices 0--19, not the original non-sequential class IDs
from the dataset paper (which go up to 34 with gaps at 2, 11, 13, 14,
16, etc.).  Our initial implementation used the non-sequential IDs,
producing completely wrong label assignments:
index 4 (which corresponds to \textit{pole} in sequential order) was
mapped as \textit{tree}; index 7 (which is \textit{sky} in the paper) was
mapped as \textit{vehicle}, and so on.

After correcting to the sequential mapping
(Table~\ref{tab:rellis_map}), training converged properly.  Sky IoU reached
$71.2\%$, tree IoU $61.1\%$, and bush IoU $20.8\%$ on the RELLIS-3D
validation set --- a drastic improvement from the completely collapsed
results obtained with the incorrect mapping.

\begin{table}[h]
\centering
\caption{Corrected RELLIS-3D sequential-index to AVMI class mapping.
         Classes with ``---'' are ignored during training.}
\label{tab:rellis_map}
\begin{tabular}{clcl}
\toprule
\textbf{Seq. idx} & \textbf{RELLIS class} & \textbf{Seq. idx} & \textbf{AVMI class} \\
\midrule
0  & void     & --- & ignore \\
1  & dirt     &  1  & Ground \\
2  & grass    & --- & ignore \\
3  & tree     &  3  & Tree   \\
4  & pole     & --- & ignore \\
5  & water    & --- & ignore \\
6  & sky      &  6  & Sky    \\
7  & vehicle  & --- & ignore \\
8  & object   & --- & ignore \\
9  & asphalt  & --- & ignore \\
10 & building & --- & ignore \\
11 & log      & --- & ignore \\
12 & person   & --- & ignore \\
13 & fence    & --- & ignore \\
14 & bush     & 14  & Bush   \\
15 & concrete & --- & ignore \\
16 & barrier  & --- & ignore \\
17 & puddle   & --- & ignore \\
18 & mud      & 18  & Ground \\
19 & rubble   & 19  & Rock   \\
\bottomrule
\end{tabular}
\end{table}

\subsection{Fine-Tuning Protocol}

All fine-tuning experiments initialise from the AVMI base model checkpoint
and train for \textbf{100,000 iterations} using:
\begin{itemize}
    \item SGD with momentum $0.9$, weight decay $4\times10^{-5}$
    \item Learning rate $0.003$ with polynomial decay (min $10^{-5}$)
    \item Linear warmup for 500 iterations
    \item Batch size: 4 (RUGD) / 2 (RELLIS, larger image resolution)
\end{itemize}
Crop size is fixed at $300\times375$ to preserve the $97\times97$ PSA
attention mask dimensions learned during AVMI scratch training.  Checkpoints
are saved every 10,000 iterations; validation mIoU is evaluated at each
checkpoint.

% ───────────────────────────────────────────────────────────────────────────
\section{Cross-Dataset Transfer Results}
\label{sec:results}

Table~\ref{tab:results} summarises validation mIoU for all experiments.
Figure~\ref{fig:ugv_qualitative} shows qualitative segmentation results on
held-out UGV images not seen during any stage of training.

\begin{table*}[t]
\centering
\caption{Validation IoU (\%) for all models after 100k fine-tuning iterations.
         The base model (row 1) is evaluated on the AVMI validation set;
         all other rows are evaluated on their respective source-dataset
         validation splits.  ``n/a'' = class not present in that mapping.
         Best result per column is \textbf{bold}.}
\label{tab:results}
\renewcommand{\arraystretch}{1.15}
\begin{tabular}{llccccccc}
\toprule
\textbf{Model} & \textbf{Strategy} &
  \textbf{Sky} & \textbf{Tree} & \textbf{Bush} &
  \textbf{Ground} & \textbf{Rock} & \textbf{Stump} & \textbf{mIoU} \\
\midrule
AVMI Base (scratch, 240k) & ---
  & 86.9 & 70.2 & 4.9 & 91.2 & 40.4 & 15.5 & 51.51 \\
\midrule
RUGD $\to$ AVMI (Full Mapping)
  & 100k FT
  & 57.0 &  5.6 & 0.0 & 88.4 & 74.3 & 68.7 & 49.0  \\
RUGD $\to$ AVMI (Selective)
  & 100k FT
  & \textbf{87.4} & 58.7 & 5.2 & \textbf{91.6} & \textbf{93.6} & n/a  & 67.4  \\
RUGD $\to$ AVMI (Full+ClassWeight)
  & 100k FT
  & 49.2 & 10.4 & 6.7 & 84.9 & 69.2 & 69.9 & 48.4  \\
\midrule
RELLIS $\to$ AVMI (Wrong index)
  & 100k FT
  & \multicolumn{6}{c}{\textit{--- training diverged (incorrect labels) ---}} & --- \\
RELLIS $\to$ AVMI (Fixed index)
  & 100k FT
  & 71.2 & \textbf{61.1} & \textbf{20.8} & 63.8 & 0.0 & 0.5 & \textbf{36.2} \\
RELLIS $\to$ AVMI (Selective)
  & 100k FT
  &  0.0 & 48.1 & 21.3 &  0.0 & n/a  & n/a  & 17.4  \\
\bottomrule
\end{tabular}
\end{table*}

\subsection{Quantitative Analysis}

The RUGD selective model achieves the highest mIoU among RUGD experiments
($67.4\%$) by learning a clean signal from only visually consistent pixels.
Tree IoU improves from $5.6\%$ (full mapping) to $58.7\%$ (selective),
confirming that the class imbalance under full mapping was the root cause
of tree detection failure.  Sky and rock achieve near-perfect IoU ($87.4\%$
and $93.6\%$ respectively) because these classes are visually distinctive
and have sufficient selective-mapping pixel counts.

The RELLIS fixed-index model achieves IoU.sky $=71.2\%$, IoU.tree $=61.1\%$,
and IoU.bush $=20.8\%$, demonstrating that RELLIS-3D is a viable transfer
source for learning sky/tree/bush segmentation once annotation indices are
correctly interpreted.  Ground IoU of $63.8\%$ is lower than for RUGD
models because RELLIS uses \textit{dirt} and \textit{mud} as ground
classes, which are less visually abundant in the dataset than RUGD's
broader ground representation.  Rock and stump IoU remain near zero because
the RELLIS dataset contains very few rubble and log pixels.

\subsection{Qualitative Evaluation on UGV Images}

\begin{figure*}[t]
    \centering
    % Row 1: RELLIS fixed-index model on UGV images
    \includegraphics[width=0.32\linewidth]{results/avmi_fixed_test_rellis_ugv/01_s1_00701.png}
    \includegraphics[width=0.32\linewidth]{results/avmi_fixed_test_rellis_ugv/02_s1_00226.png}
    \includegraphics[width=0.32\linewidth]{results/avmi_fixed_test_rellis_ugv/03_s1_00480.png}
    \\[4pt]
    % Row 2: RUGD selective model on UGV images
    \includegraphics[width=0.32\linewidth]{results/avmi_selective_test_rugd_ugv/01_s1_01783.png}
    \includegraphics[width=0.32\linewidth]{results/avmi_selective_test_rugd_ugv/02_s1_04144.png}
    \includegraphics[width=0.32\linewidth]{results/avmi_selective_test_rugd_ugv/03_s2_01281.png}
    \caption{Qualitative segmentation on held-out UGV images (not seen during
             training).
             \textbf{Top row}: RELLIS fixed-index fine-tuned model.
             Sky (blue) and tree (green) are correctly segmented; ground
             (brown) covers the open grass area.
             \textbf{Bottom row}: RUGD selective fine-tuned model.
             Sky and tree detection is strong within the RUGD domain, but
             outputs are less stable on UGV images due to the 94\% pixel
             ignore rate during training.
             Each panel shows (left) original image, (centre) segmentation
             mask, (right) overlay with per-class percentages.}
    \label{fig:ugv_qualitative}
\end{figure*}

Figure~\ref{fig:ugv_qualitative} compares the best RELLIS and RUGD models
on unseen UGV images.  The RELLIS fixed-index model (top row) produces
coherent segmentation: sky occupies the upper portion of the frame
(18--29\% of pixels), tree canopy is correctly identified along frame edges
and depth (15--30\%), and the grassy ground is the dominant traversable
region (47--58\%).  The RUGD selective model (bottom row) detects sky and
tree accurately within RUGD-domain imagery but produces less stable outputs
on UGV images because the high ignore rate (94\% of pixels) left it without
supervision for the bright-green grass texture prevalent in the UGV dataset.

Figure~\ref{fig:cross_dataset} shows the base model applied directly to
RUGD and RELLIS images without any fine-tuning, demonstrating that the
AVMI-pretrained features already capture tree and sky structure, validating
the transfer learning premise.

\begin{figure}[h]
    \centering
    \includegraphics[width=0.49\linewidth]{results/cross_dataset_test/rugd_1_trail-12_01011.png}
    \includegraphics[width=0.49\linewidth]{results/cross_dataset_test/rellis_1_frame003906-1581797541_009.png}
    \caption{AVMI base model (no fine-tuning) applied to RUGD (left) and
             RELLIS-3D (right) images.  Despite being trained only on
             synthetic UGV data, the model correctly identifies sky and
             tree regions in both real datasets, confirming that the
             learned feature representations generalise across domains.}
    \label{fig:cross_dataset}
\end{figure}

\subsection{Summary of Findings}

Our experiments surface three key lessons for cross-dataset transfer in
off-road segmentation:

\begin{enumerate}
    \item \textbf{Annotation format must be verified end-to-end.}
          The RELLIS-3D sequential index error caused complete training
          failure that would have been invisible without per-class IoU
          monitoring.  A $>$50\% absolute improvement in tree IoU was
          recovered by a single lookup-table correction.

    \item \textbf{Class imbalance kills rare classes.}
          Under full RUGD mapping, ground (68\% of pixels) suppresses
          tree (0.2\%) regardless of architecture or learning rate.
          Selective mapping is the most effective mitigation for source
          datasets with highly skewed distributions, at the cost of reduced
          coverage.

    \item \textbf{RELLIS-3D is the stronger transfer source for UGV
          terrain segmentation.}
          RELLIS forest trail imagery shares more visual structure with
          UGV navigation environments than RUGD, and its naturally balanced
          class distribution (sky/grass/tree/bush = 94\% of pixels)
          aligns well with the visual classes that matter most for UGV
          perception.
\end{enumerate}

% ───────────────────────────────────────────────────────────────────────────
%  REFERENCES  (add to your main bibliography)
% ───────────────────────────────────────────────────────────────────────────
%
% \bibitem{rugd2019}
%   O. Wigness et al., "A RUGD Dataset for Autonomous Navigation and Visual
%   Perception in Unstructured Outdoor Environments," IROS 2019.
%
% \bibitem{rellis2021}
%   P. Jiang et al., "RELLIS-3D Dataset: Data, Benchmarks and Analysis,"
%   ICRA 2021.
%
% \bibitem{ganav2022}
%   H. Gao et al., "GANav: Group-wise Attention Network for Classifying
%   Navigable Regions in Unstructured Outdoor Environments," IEEE RAL 2022.
